% Part: computability
% Chapter: computability-theory
% Section: s-m-n

\documentclass[../../../include/open-logic-section]{subfiles}

\begin{document}

\olfileid{cmp}{thy}{smn}
\olsection{$s$-$m$-$n$ دا قضیہ}

\begin{explain}
اگلا قضیہ ``$s$-$m$-$n$ دا قضیہ'' کہلاندا اے؛ ایس ناں دی
وجہ ہُنّے صاف ہو جائے گی۔ اوکھا کم ایہ سمجھنا اے کہ
قضیہ ٹھیک ٹھیک کیہ کہندا اے؛ جے بیان سمجھ آ جائے، تاں
نتیجہ خاصا ظاہر لگے گا۔
\end{explain}

\begin{thm}
\ollabel{thm:s-m-n}
  قدرتی عدداں دے ہر جوڑے $n$ تے~$m$ لئی اک بنیادی
  بازگشتی فنکشن~$s^m_n$ موجود اے، ایداں کہ ہر سلسلے
  $e$، $a_0$، \dots، $a_{m-1}$، $y_0$، \dots، $y_{n-1}$ لئی
  \[
  \cfind{s^m_n(e, a_0, \dots, a_{m-1})}[n](y_0, \dots, y_{n-1}) \simeq
  \cfind{e}[m+n](a_0, \dots, a_{m-1}, y_0, \dots, y_{n-1}).
\]
\end{thm}

\begin{explain}
% PNB-CMP-011: source B7 دے دو explanatory x پروگرام دے اشاریہ e نال بدلے؛ theorem unchanged۔
$s^m_n$ نوں \emph{پروگراماں} اُتے کم کرن والا سمجھنا مددگار
اے۔ یعنی $s^m_n$ اک پروگرام~$e$ لیندا اے، جیہڑا $(m+n)$
آرگومنٹاں والے فنکشن لئی اے، تے نال مقرر انپٹ $a_0$،
\dots، $a_{m-1}$ وی؛ تے اوہ پروگرام
$s^m_n(e, a_0, \dots, a_{m-1})$ واپس دیندا اے، جیہڑا باقی
آرگومنٹاں دے $n$-آرگومنٹ والے فنکشن لئی اے۔ \iftag{TMs}{جے
تسی $e$ نوں اک ٹیورنگ مشین دا بیان سمجھو، تاں
$s^m_n(e, a_0, \dots, a_{m-1})$ اوہ ٹیورنگ مشین اے جیہڑی
انپٹ $y_0$، \dots،~$y_{n-1}$ ملّن اُتے، $a_0$، \dots،~$a_{m-1}$
نوں انپٹ لڑی دے مُڈھ وچ لا دیندی اے، تے~$e$ چلا دیندی
اے۔ فیر ہر $s^m_n$ صرف اک بنیادی بازگشتی فنکشن اے جیہڑا
مناسب ٹیورنگ مشین دا کوڈ لبھدا اے۔}{}
\end{explain}

\end{document}
